\documentclass[11pt,a4paper]{article}
\usepackage[margin=1in]{geometry}
\usepackage{amsmath,amssymb,amsthm}
\usepackage{graphicx}
\usepackage{microtype}
\usepackage[hidelinks]{hyperref}

\newtheorem{theorem}{Theorem}
\newtheorem{lemma}[theorem]{Lemma}
\newtheorem{corollary}[theorem]{Corollary}
\newcommand{\R}{\mathbb{R}}
\newcommand{\VLIP}{\mathsf{VL}_{\mathsf{IP}}}

\title{Atomic direct limits for interval-preserving maps}
\author{A partial answer to an open question in arXiv:2208.13813}
\date{17 August 2026}

\begin{document}
\maketitle

\noindent\textbf{Status.}
Candidate substantial partial result, likely valid, pending human review.
The general existence question remains open, but every direct system in
$\VLIP$ whose objects are lattice-isomorphic to spaces $c_{00}(A)$ has a
direct limit in $\VLIP$.  This includes systems of finite-dimensional
\emph{Archimedean} vector lattices.  The qualifier is necessary because the
source allows non-Archimedean vector lattices.

\section{Source question}

Ding and de Jeu study direct limits in categories of vector lattices, normed
lattices, and Banach lattices~\cite{DingDeJeu}.  On source PDF page~25 they
state that it is open whether direct limits always exist in each of six
exceptional categories.  The first is $\VLIP$: real vector lattices with
interval-preserving linear maps.  The other five impose norm or completeness
conditions.  The theorem below solves a broad algebraically atomic subcase of
the first question.

\begin{center}
\includegraphics[width=\textwidth]{figures/open_problem_crop.png}
\end{center}

A positive linear map $T:E\to F$ is interval preserving when
\[
 T([0,x]_E)=[0,Tx]_F\qquad(x\in E_+).
\]
An atom is a nonzero $a\geq0$ for which $[0,a]$ is contained in the
one-dimensional ray $\R a$.  For a set $A$, $c_{00}(A)$ denotes the vector
lattice of finitely supported real functions on $A$, with coordinatewise
order.

\section{Proof intuition}

An interval-preserving map out of $c_{00}(A)$ has an elementary form: every
coordinate atom is sent either to zero or to a positive multiple of an atom.
Different coordinates may merge onto the same atomic ray.  Thus a direct
system is combinatorially a directed system of weighted partial maps between
sets of atoms.  In its ordinary vector-space direct limit, the surviving
atomic rays form a basis.  Giving that basis the coordinatewise order produces
the desired vector lattice.  The same atomic description forces every
factorization map in the universal property to be interval preserving.

\section{Atomic characterization}

\begin{lemma}\label{lem:monomial}
Let $A$ be a set, let $F$ be a vector lattice, and let
$T:c_{00}(A)\to F$ be linear.  Then $T$ is interval preserving if and only if
for every coordinate vector $e_a$, the vector $Te_a$ is either zero or an
atom of $F$.
\end{lemma}

\begin{proof}
If $T$ is interval preserving, it is positive and
\[
 [0,Te_a]=T([0,e_a])=\{tTe_a:0\leq t\leq1\}.
\]
Consequently a nonzero $Te_a$ is an atom.

Conversely, suppose that every $Te_a$ is zero or an atom.  Then $T$ is
positive.  Fix $x=\sum_{a\in S}x_ae_a\geq0$, where $S$ is finite, and let
$0\leq y\leq Tx$.  Group the nonzero vectors $Te_a$ for $a\in S$ according
to their positive rays.  Choose one atom $b_r$ on every resulting ray and
write
\[
 Tx=\sum_r s_rb_r,
\]
where $s_r>0$ and the $b_r$ are pairwise disjoint.  The Riesz decomposition
property gives $y=\sum_r y_r$ with $0\leq y_r\leq s_rb_r$.  Atomicity gives
$y_r=t_rb_r$ for some $0\leq t_r\leq s_r$.

Within a fixed group write $Te_a=\beta_ab_r$.  Since
$s_r=\sum_a x_a\beta_a$, one can choose numbers
$0\leq z_a\leq x_a$ with $\sum_a z_a\beta_a=t_r$; this is just the fact that
the sum of the intervals $[0,x_a\beta_a]$ is $[0,s_r]$.  Doing this for every
group and putting $z=\sum_{a\in S}z_ae_a$ gives $0\leq z\leq x$ and
$Tz=y$.  Hence $[0,Tx]\subseteq T([0,x])$; the reverse inclusion follows
from positivity.
\end{proof}

\section{Existence of the direct limit}

\begin{theorem}\label{thm:limit}
Let $(E_i,\phi_{ji})$ be a direct system over an arbitrary directed set in
$\VLIP$.  Suppose that every $E_i$ is lattice-isomorphic to $c_{00}(A_i)$
for some set $A_i$.  Then the system has a direct limit in $\VLIP$, and its
underlying vector lattice is lattice-isomorphic to $c_{00}(C)$ for a set
$C$ of surviving atomic rays.
\end{theorem}

\begin{proof}
Choose atomic coordinate vectors $e_{i,a}$ in every $E_i$.  By
Lemma~\ref{lem:monomial}, each connecting map sends $e_{i,a}$ either to zero
or to a positive multiple of a coordinate atom at the later stage.

Let $(V,q_i)$ be the ordinary algebraic direct limit of the underlying vector
spaces.  Consider all nonzero vectors $q_i(e_{i,a})$ and identify two of them
when they lie on the same positive ray in $V$.  Let $C$ be the set of the
resulting rays, and choose a positive generator $u_c$ from every ray.

We claim that $(u_c)_{c\in C}$ is a vector-space basis of $V$.  It spans:
every element of $V$ is the image of an element at one stage, and every such
element has finite atomic support.  For independence, take a finite relation
among distinct $u_c$ and represent all of its terms by coordinate atoms at
their respective stages.  Directedness provides one later stage $k$ receiving
all of them.  No representative can map to zero there, since its image in
$V$ is nonzero.  Each therefore maps to a positive multiple of a coordinate
atom of $E_k$.  Two representatives cannot land on the same coordinate atom,
since then their images in $V$ would lie on the same positive ray.  Their
images at stage $k$ are consequently distinct coordinate vectors and are
linearly independent.  The original relation is trivial.

Use this basis to order $V$ coordinatewise:
\[
 V_+=\left\{\sum_{c\in C}\alpha_cu_c:
        \alpha_c\geq0\text{ and only finitely many }\alpha_c\ne0\right\}.
\]
Then $V$ is a vector lattice isomorphic to $c_{00}(C)$.  Each $q_i$ sends a
coordinate atom to zero or to a positive multiple of a coordinate atom of
$V$, so Lemma~\ref{lem:monomial} says that every $q_i$ is interval preserving.

It remains to check the universal property in the smaller category, not only
among vector spaces.  Let $T_i:E_i\to F$ be any compatible cocone of
interval-preserving maps into a vector lattice $F$.  The vector-space direct
limit supplies a unique linear map $T:V\to F$ with $Tq_i=T_i$.  Every basis
atom $u_c$ is a positive multiple of some $q_i(e_{i,a})$, so $Tu_c$ is a
positive multiple of $T_i(e_{i,a})$.  Lemma~\ref{lem:monomial} shows that
$Tu_c$ is zero or an atom of $F$.  Applying that lemma once more proves that
$T$ is interval preserving.  Thus $(V,q_i)$ has precisely the universal
property of a direct limit in $\VLIP$.
\end{proof}

\begin{corollary}
Every direct system of finite-dimensional Archimedean vector lattices and
interval-preserving linear maps has a direct limit in $\VLIP$.
\end{corollary}

\begin{proof}
Every finite-dimensional Archimedean vector lattice is lattice-isomorphic to
$\R^n=c_{00}(\{1,\ldots,n\})$ with coordinatewise order, so
Theorem~\ref{thm:limit} applies.
\end{proof}

\section{Upgrade attempts, novelty, and limits}

The first proof attempt covered only coordinatewise finite-dimensional
objects.  The key observation that every element uses finitely many atoms
upgrades the argument to arbitrary $c_{00}(A_i)$ and arbitrary directed index
sets.  A further attempt to remove Archimedeanity from the finite-dimensional
corollary does not follow from this method: finite-dimensional
non-Archimedean lattices, such as lexicographically ordered spaces, need not be
spanned by atoms.  Likewise $\ell^\infty(A)$ is atomic in the order-theoretic
sense but is not the algebraic span of its atoms, so the basis construction
does not cover it.

The result also does not automatically extend to the five normed or Banach
categories in the source question.  Quotient norms and completion introduce
conditions absent from the algebraic colimit, and completion can add elements
that are not finite sums of atoms; interval preservation need not survive the
last inclusion into the completion.

A bounded novelty search covered the run-wide result and attempt indexes, the
arXiv identifier, exact phrases around the source question, and combinations
of ``atomic vector lattice'', ``interval preserving'', and ``direct limit''.
No primary source stating this atomic subcase was found.  Because the proof is
elementary and categorical, novelty confidence is moderate; correctness
confidence is high.  Human review should focus on the basis-independence
argument and the repeated use of Lemma~\ref{lem:monomial} in the universal
property.

\begin{thebibliography}{1}
\bibitem{DingDeJeu}
C.~Ding and M.~de Jeu,
\emph{Direct limits in categories of normed vector lattices and Banach
lattices}, arXiv:2208.13813.
\end{thebibliography}

\end{document}
